% --- Archon physics formalization source begin ---
% archon:physics
% archon:covers PhyXMiniProblems/problem_phyx_mini_0102.lean
% archon:source-report reports/phyx_mini/problem_phyx_mini_0102.source.json

\chapter{Physics problem phyx\_mini\_0102}
\label{ch:PhyXMiniProblems_problem_phyx_mini_0102}

\paragraph{Problem source.}
\#\# Physical scenario

Light is incident in air at an angle \$\textbackslash{}theta\_a\$ (figure) on the upper surface of a transparent plate, the surfaces of the plate being plane and parallel to each other.A ray of light is incident at an angle of \$66.0\textasciicircum{}\textbackslash{}circ\$ on one surface of a glass plate \$2.40 \textbackslash{}, \textbackslash{}text\{cm\}\$ thick with an index of refraction of \$1.80\$. The medium on either side of the plate is air.

\#\# Question

Find the lateral displacement between the incident and emergent rays.

\#\# Answer choices

- A: A: 1.62cm

- B: B: 1.54cm

- C: C: 1.75cm

- D: D: 1.59cm

\#\# Figure caption (auxiliary; use the image as primary evidence)

The image is a diagram depicting the path of light through a rectangular medium.

- **Objects and Elements**:
  - A rectangular region with a light blue fill representing a medium with refractive index \textbackslash{}( n' \textbackslash{}).
  - The surrounding area has a refractive index \textbackslash{}( n \textbackslash{}).
  - A purple line with arrows represents the light path, which enters, refracts through, and exits the medium.

- **Light Path**:
  - The light enters the medium at point \textbackslash{}( P \textbackslash{}) with an angle of incidence \textbackslash{}( \textbackslash{}theta\_a \textbackslash{}).
  - Inside the medium, there's an angle of refraction \textbackslash{}( \textbackslash{}theta\_b \textbackslash{}) as the light travels through.
  - The light exits the medium at point \textbackslash{}( Q \textbackslash{}) with an angle of exit \textbackslash{}( \textbackslash{}theta'\_a \textbackslash{}).

- **Angles and Distances**:
  - The normal lines at entry and exit points are indicated with dashed lines.
  - Angles are marked as \textbackslash{}( \textbackslash{}theta\_a \textbackslash{}), \textbackslash{}( \textbackslash{}theta'\_b \textbackslash{}), \textbackslash{}( \textbackslash{}theta\_b \textbackslash{}), and \textbackslash{}( \textbackslash{}theta'\_a \textbackslash{}).
  - The height \textbackslash{}( t \textbackslash{}) of the medium and the lateral displacement \textbackslash{}( d \textbackslash{}) of the light path are indicated with double-headed arrows.

- **Relationships**:
  - The relationship between the angles of incidence and refraction and between different media is depicted, illustrating principles of light refraction.

- **Other Notations**:
  - The diagram is likely explaining refraction or Snell's Law, demonstrating how light bends as it moves through different media.

\#\# Dataset metadata

Geometrical Optics; Spatial Relation Reasoning, Physical Model Grounding Reasoning

\paragraph{Recorded answer/context.}
A: A: 1.62cm

\paragraph{Figure/image path.}
/root/proposal\_for\_physic/science-mango/phyx\_mini\_run/phyx\_data/test\_image/102.png

\paragraph{Formalization target.}
create a compiling Lean file with sorry bodies at `PhyXMiniProblems/problem\_phyx\_mini\_0102.lean`. The Lean declarations must preserve the physical quantities, dimensions or dimensional roles, figure labels, governing-law hypotheses, and final relation expressed by this problem.
Use Mathlib/Physlib names found through LeanExplore where available. If a physics API is missing, introduce faithful local abstractions rather than scalar placeholder aliases.

\begin{theorem}[Physics formalization target]
\label{thm:physics:phyx_mini_0102:target}
\lean{PhyXMiniProblems.ProblemPhyXMini0102.lateralDisplacementIsAnswerA}
\uses{def:physics:phyx-mini-0102:phyxminiproblems-problemphyxmini0102-parallelglassplatesetup, def:physics:phyx-mini-0102:phyxminiproblems-problemphyxmini0102-matchesproblemdata, def:physics:phyx-mini-0102:phyxminiproblems-problemphyxmini0102-hasphysicalopticalparameters, def:physics:phyx-mini-0102:phyxminiproblems-problemphyxmini0102-satisfiesparallelplategeometry, def:physics:phyx-mini-0102:phyxminiproblems-problemphyxmini0102-obeyssnellslaw, def:physics:phyx-mini-0102:phyxminiproblems-problemphyxmini0102-matchesanswertonearesthundredthcentimeter}
The assigned autoformalize agent should translate this physics problem into Lean declarations in the covered file, with theorem and lemma proof bodies written as `by sorry`.
\end{theorem}
\begin{proof}
This is an autoformalization task, not a proof task. Produce statements that can later be proved by the physics prover without weakening the physical model.
\end{proof}

% --- Archon declaration topology begin ---
\section{Declaration topology}
\begin{definition}[Dim Length]
  \label{def:physics:phyx-mini-0102:phyxminiproblems-problemphyxmini0102-dimlength}
  \lean{PhyXMiniProblems.ProblemPhyXMini0102.DimLength}
  A physical length, independent of the unit system used to read it
\end{definition}
\begin{proof}
  This is the declared model definition of Dim Length.
\end{proof}
\begin{definition}[centimeter Unit Choices]
  \label{def:physics:phyx-mini-0102:phyxminiproblems-problemphyxmini0102-centimeterunitchoices}
  \lean{PhyXMiniProblems.ProblemPhyXMini0102.centimeterUnitChoices}
  Unit choices in which physical lengths are read in centimeters
\end{definition}
\begin{proof}
  This is the declared model definition of centimeter Unit Choices.
\end{proof}
\begin{definition}[value In Centimeters]
  \label{def:physics:phyx-mini-0102:phyxminiproblems-problemphyxmini0102-valueincentimeters}
  \lean{PhyXMiniProblems.ProblemPhyXMini0102.valueInCentimeters}
  \uses{def:physics:phyx-mini-0102:phyxminiproblems-problemphyxmini0102-dimlength, def:physics:phyx-mini-0102:phyxminiproblems-problemphyxmini0102-centimeterunitchoices}
  The real-valued centimeter readout of a dimensionful length
\end{definition}
\begin{proof}
  This is the declared model definition of value In Centimeters.
\end{proof}
\begin{definition}[degrees To Radians]
  \label{def:physics:phyx-mini-0102:phyxminiproblems-problemphyxmini0102-degreestoradians}
  \lean{PhyXMiniProblems.ProblemPhyXMini0102.degreesToRadians}
  Convert a degree readout to the radian scalar used by Real.sin
\end{definition}
\begin{proof}
  This is the declared model definition of degrees To Radians.
\end{proof}
\begin{definition}[Diagram Plane]
  \label{def:physics:phyx-mini-0102:phyxminiproblems-problemphyxmini0102-diagramplane}
  \lean{PhyXMiniProblems.ProblemPhyXMini0102.DiagramPlane}
  The two-dimensional plane of the ray diagram, with centimeter coordinates
\end{definition}
\begin{proof}
  This is the declared model definition of Diagram Plane.
\end{proof}
\begin{definition}[Ray Direction]
  \label{def:physics:phyx-mini-0102:phyxminiproblems-problemphyxmini0102-raydirection}
  \lean{PhyXMiniProblems.ProblemPhyXMini0102.RayDirection}
  \uses{def:physics:phyx-mini-0102:phyxminiproblems-problemphyxmini0102-diagramplane}
  A propagation direction is represented by a nonzero diagram vector
\end{definition}
\begin{proof}
  This is the declared model definition of Ray Direction.
\end{proof}
\begin{definition}[Optical Region]
  \label{def:physics:phyx-mini-0102:phyxminiproblems-problemphyxmini0102-opticalregion}
  \lean{PhyXMiniProblems.ProblemPhyXMini0102.OpticalRegion}
  The three homogeneous optical regions traversed by the ray
\end{definition}
\begin{proof}
  This is the declared model definition of Optical Region.
\end{proof}
\begin{definition}[Plate Face]
  \label{def:physics:phyx-mini-0102:phyxminiproblems-problemphyxmini0102-plateface}
  \lean{PhyXMiniProblems.ProblemPhyXMini0102.PlateFace}
  The plane and parallel upper and lower faces of the plate
\end{definition}
\begin{proof}
  This is the declared model definition of Plate Face.
\end{proof}
\begin{definition}[Ray Segment]
  \label{def:physics:phyx-mini-0102:phyxminiproblems-problemphyxmini0102-raysegment}
  \lean{PhyXMiniProblems.ProblemPhyXMini0102.RaySegment}
  The three directed portions of the light path
\end{definition}
\begin{proof}
  This is the declared model definition of Ray Segment.
\end{proof}
\begin{definition}[Figure Point]
  \label{def:physics:phyx-mini-0102:phyxminiproblems-problemphyxmini0102-figurepoint}
  \lean{PhyXMiniProblems.ProblemPhyXMini0102.FigurePoint}
  Point labels retained from the primary figure
\end{definition}
\begin{proof}
  This is the declared model definition of Figure Point.
\end{proof}
\begin{definition}[incident Region]
  \label{def:physics:phyx-mini-0102:phyxminiproblems-problemphyxmini0102-incidentregion}
  \lean{PhyXMiniProblems.ProblemPhyXMini0102.incidentRegion}
  \uses{def:physics:phyx-mini-0102:phyxminiproblems-problemphyxmini0102-opticalregion, def:physics:phyx-mini-0102:phyxminiproblems-problemphyxmini0102-plateface}
  Incident optical region at each face, in propagation order
\end{definition}
\begin{proof}
  This is the declared model definition of incident Region.
\end{proof}
\begin{definition}[transmitted Region]
  \label{def:physics:phyx-mini-0102:phyxminiproblems-problemphyxmini0102-transmittedregion}
  \lean{PhyXMiniProblems.ProblemPhyXMini0102.transmittedRegion}
  \uses{def:physics:phyx-mini-0102:phyxminiproblems-problemphyxmini0102-opticalregion, def:physics:phyx-mini-0102:phyxminiproblems-problemphyxmini0102-plateface}
  Transmitted optical region at each face, in propagation order
\end{definition}
\begin{proof}
  This is the declared model definition of transmitted Region.
\end{proof}
\begin{definition}[incident Segment]
  \label{def:physics:phyx-mini-0102:phyxminiproblems-problemphyxmini0102-incidentsegment}
  \lean{PhyXMiniProblems.ProblemPhyXMini0102.incidentSegment}
  \uses{def:physics:phyx-mini-0102:phyxminiproblems-problemphyxmini0102-plateface, def:physics:phyx-mini-0102:phyxminiproblems-problemphyxmini0102-raysegment}
  Ray segment incident on each plate face
\end{definition}
\begin{proof}
  This is the declared model definition of incident Segment.
\end{proof}
\begin{definition}[transmitted Segment]
  \label{def:physics:phyx-mini-0102:phyxminiproblems-problemphyxmini0102-transmittedsegment}
  \lean{PhyXMiniProblems.ProblemPhyXMini0102.transmittedSegment}
  \uses{def:physics:phyx-mini-0102:phyxminiproblems-problemphyxmini0102-plateface, def:physics:phyx-mini-0102:phyxminiproblems-problemphyxmini0102-raysegment}
  Ray segment transmitted through each plate face
\end{definition}
\begin{proof}
  This is the declared model definition of transmitted Segment.
\end{proof}
\begin{definition}[Parallel Glass Plate Setup]
  \label{def:physics:phyx-mini-0102:phyxminiproblems-problemphyxmini0102-parallelglassplatesetup}
  \lean{PhyXMiniProblems.ProblemPhyXMini0102.ParallelGlassPlateSetup}
  \uses{def:physics:phyx-mini-0102:phyxminiproblems-problemphyxmini0102-dimlength, def:physics:phyx-mini-0102:phyxminiproblems-problemphyxmini0102-diagramplane, def:physics:phyx-mini-0102:phyxminiproblems-problemphyxmini0102-raydirection, def:physics:phyx-mini-0102:phyxminiproblems-problemphyxmini0102-opticalregion, def:physics:phyx-mini-0102:phyxminiproblems-problemphyxmini0102-plateface, def:physics:phyx-mini-0102:phyxminiproblems-problemphyxmini0102-raysegment, def:physics:phyx-mini-0102:phyxminiproblems-problemphyxmini0102-figurepoint}
  Physical quantities and figure readouts for the parallel glass plate. For a face, incidenceAngleRadians is the angle on the incident side and refractionAngleRadians is the angle on the transmitted side. Consequently the four values carry the figure labels θ\_a, θ'\_b, θ\_b, and θ'\_a at the entry-incidence, entry-refraction, exit-incidence, and exit-refraction positions respectively
\end{definition}
\begin{proof}
  This is the declared model definition of Parallel Glass Plate Setup.
\end{proof}
\begin{definition}[angle Between Directions]
  \label{def:physics:phyx-mini-0102:phyxminiproblems-problemphyxmini0102-anglebetweendirections}
  \lean{PhyXMiniProblems.ProblemPhyXMini0102.angleBetweenDirections}
  \uses{def:physics:phyx-mini-0102:phyxminiproblems-problemphyxmini0102-raydirection}
  The undirected radian angle between two nonzero diagram directions
\end{definition}
\begin{proof}
  This is the declared model definition of angle Between Directions.
\end{proof}
\begin{definition}[Is Physical Acute Angle]
  \label{def:physics:phyx-mini-0102:phyxminiproblems-problemphyxmini0102-isphysicalacuteangle}
  \lean{PhyXMiniProblems.ProblemPhyXMini0102.IsPhysicalAcuteAngle}
  An angle readout lies on the physical acute branch of geometrical optics
\end{definition}
\begin{proof}
  This is the declared model definition of Is Physical Acute Angle.
\end{proof}
\begin{definition}[Matches Problem Data]
  \label{def:physics:phyx-mini-0102:phyxminiproblems-problemphyxmini0102-matchesproblemdata}
  \lean{PhyXMiniProblems.ProblemPhyXMini0102.MatchesProblemData}
  \uses{def:physics:phyx-mini-0102:phyxminiproblems-problemphyxmini0102-valueincentimeters, def:physics:phyx-mini-0102:phyxminiproblems-problemphyxmini0102-degreestoradians, def:physics:phyx-mini-0102:phyxminiproblems-problemphyxmini0102-parallelglassplatesetup}
  Problem-statement readouts: a transparent 2.40 cm glass plate of refractive index 1.80, a 66.0° incidence angle in air, and air of index one on both sides. No value of the requested lateral displacement occurs here
\end{definition}
\begin{proof}
  This is the declared model definition of Matches Problem Data.
\end{proof}
\begin{definition}[Has Physical Optical Parameters]
  \label{def:physics:phyx-mini-0102:phyxminiproblems-problemphyxmini0102-hasphysicalopticalparameters}
  \lean{PhyXMiniProblems.ProblemPhyXMini0102.HasPhysicalOpticalParameters}
  \uses{def:physics:phyx-mini-0102:phyxminiproblems-problemphyxmini0102-valueincentimeters, def:physics:phyx-mini-0102:phyxminiproblems-problemphyxmini0102-parallelglassplatesetup, def:physics:phyx-mini-0102:phyxminiproblems-problemphyxmini0102-isphysicalacuteangle}
  Positivity and principal-angle-branch conditions for the optical model
\end{definition}
\begin{proof}
  This is the declared model definition of Has Physical Optical Parameters.
\end{proof}
\begin{definition}[Satisfies Parallel Plate Geometry]
  \label{def:physics:phyx-mini-0102:phyxminiproblems-problemphyxmini0102-satisfiesparallelplategeometry}
  \lean{PhyXMiniProblems.ProblemPhyXMini0102.SatisfiesParallelPlateGeometry}
  \uses{def:physics:phyx-mini-0102:phyxminiproblems-problemphyxmini0102-valueincentimeters, def:physics:phyx-mini-0102:phyxminiproblems-problemphyxmini0102-incidentsegment, def:physics:phyx-mini-0102:phyxminiproblems-problemphyxmini0102-transmittedsegment, def:physics:phyx-mini-0102:phyxminiproblems-problemphyxmini0102-parallelglassplatesetup, def:physics:phyx-mini-0102:phyxminiproblems-problemphyxmini0102-anglebetweendirections}
  The plane-parallel surface geometry and the P--Q construction shown in the figure. The final field is the general parallel-plate displacement relation d = t sin (θ\_a - θ'\_b) / cos θ'\_b. It relates arbitrary physical readouts and does not assign the requested numerical value of d
\end{definition}
\begin{proof}
  This is the declared model definition of Satisfies Parallel Plate Geometry.
\end{proof}
\begin{definition}[Satisfies Snells Law At]
  \label{def:physics:phyx-mini-0102:phyxminiproblems-problemphyxmini0102-satisfiessnellslawat}
  \lean{PhyXMiniProblems.ProblemPhyXMini0102.SatisfiesSnellsLawAt}
  \uses{def:physics:phyx-mini-0102:phyxminiproblems-problemphyxmini0102-plateface, def:physics:phyx-mini-0102:phyxminiproblems-problemphyxmini0102-incidentregion, def:physics:phyx-mini-0102:phyxminiproblems-problemphyxmini0102-transmittedregion, def:physics:phyx-mini-0102:phyxminiproblems-problemphyxmini0102-parallelglassplatesetup}
  Snell's law n₁ sin θ₁ = n₂ sin θ₂ at one plate face
\end{definition}
\begin{proof}
  This is the declared model definition of Satisfies Snells Law At.
\end{proof}
\begin{definition}[Obeys Snells Law]
  \label{def:physics:phyx-mini-0102:phyxminiproblems-problemphyxmini0102-obeyssnellslaw}
  \lean{PhyXMiniProblems.ProblemPhyXMini0102.ObeysSnellsLaw}
  \uses{def:physics:phyx-mini-0102:phyxminiproblems-problemphyxmini0102-parallelglassplatesetup, def:physics:phyx-mini-0102:phyxminiproblems-problemphyxmini0102-satisfiessnellslawat}
  The ray obeys Snell's law at the entry and exit faces
\end{definition}
\begin{proof}
  This is the declared model definition of Obeys Snells Law.
\end{proof}
\begin{definition}[Answer Choice]
  \label{def:physics:phyx-mini-0102:phyxminiproblems-problemphyxmini0102-answerchoice}
  \lean{PhyXMiniProblems.ProblemPhyXMini0102.AnswerChoice}
  The four displayed answer labels
\end{definition}
\begin{proof}
  This is the declared model definition of Answer Choice.
\end{proof}
\begin{definition}[displacement Centimeters]
  \label{def:physics:phyx-mini-0102:phyxminiproblems-problemphyxmini0102-answerchoice-displacementcentimeters}
  \lean{PhyXMiniProblems.ProblemPhyXMini0102.AnswerChoice.displacementCentimeters}
  \uses{def:physics:phyx-mini-0102:phyxminiproblems-problemphyxmini0102-answerchoice}
  Lateral-displacement readout in centimeters printed beside each choice
\end{definition}
\begin{proof}
  This is the declared model definition of displacement Centimeters.
\end{proof}
\begin{definition}[recorded Answer Choice]
  \label{def:physics:phyx-mini-0102:phyxminiproblems-problemphyxmini0102-recordedanswerchoice}
  \lean{PhyXMiniProblems.ProblemPhyXMini0102.recordedAnswerChoice}
  \uses{def:physics:phyx-mini-0102:phyxminiproblems-problemphyxmini0102-answerchoice}
  Dataset metadata recording answer A; this definition is not a premise
\end{definition}
\begin{proof}
  This is the declared model definition of recorded Answer Choice.
\end{proof}
\begin{definition}[Matches Answer To Nearest Hundredth Centimeter]
  \label{def:physics:phyx-mini-0102:phyxminiproblems-problemphyxmini0102-matchesanswertonearesthundredthcentimeter}
  \lean{PhyXMiniProblems.ProblemPhyXMini0102.MatchesAnswerToNearestHundredthCentimeter}
  \uses{def:physics:phyx-mini-0102:phyxminiproblems-problemphyxmini0102-dimlength, def:physics:phyx-mini-0102:phyxminiproblems-problemphyxmini0102-valueincentimeters, def:physics:phyx-mini-0102:phyxminiproblems-problemphyxmini0102-answerchoice}
  A physical length rounds to a displayed hundredth-centimeter answer when its centimeter readout differs from that answer by at most 0.005 cm
\end{definition}
\begin{proof}
  This is the declared model definition of Matches Answer To Nearest Hundredth Centimeter.
\end{proof}
% --- Archon declaration topology end ---
% --- Archon physics formalization source end ---
